% ============================================================
% 分片：s02e  对应原文 PDF 第 18--21 页（并含第 22 页首句，见文末跨页注）
% 负责章节：§2.3 Twisting the worldsheet；2.3.1；2.3.2；2.3.3
% 译者：translator-03   审校状态：初稿
% ============================================================

% ==== 原文 PDF 第 18 页 ====
\subsection{扭变世界面}

现在转向与弦论明显相关的二维 $\mathcal{N}=(2,2)$ 量子场论的一个特例：中心荷 $c=3d$ 的超共形场论。对超共形 $\sigma$ 模型，$d$ 表示靶空间的维数，因此 $c=9$ 可以描述一个 Calabi--Yau 三维流形。由于矢量与轴 $R$ 对称性都被保持，代数中不允许出现中心项。虽然镜像对称是物理弦论的对偶，但对理论作拓扑扭变将使我们能在镜像对偶的两侧提取可计算的可观测量。

\subsubsection{手征环与扭手征环}

如上所见，拓扑超荷 $Q_A$ 与 $Q_B$ 在镜像映射下互换。对超共形理论，两个 $U(1)$ $R$ 对称性都被保持，故 $Q_A$ 与 $Q_B$ 都给出 $\mathbb{Z}$-分次的拓扑扭变。可以对这两种超荷中的任意一个考虑态或算子的 $Q$ 上同调（$Q=Q_A$ 或 $Q_B$）——在最受关注的情形，即理论超共形时，态与算子之间当然存在同构。本节以及全篇讲义的关注点是算子视角。

局域算子 $O$ 若满足 $[Q_B,O]=0$，称为手征的\footnote{本节全程沿用 [9] 的术语与记号。}；若满足 $[Q_A,O]=0$，称为扭手征的。用第 2.1 节的语言：手征算子是 $Q_B$-闭的，扭手征算子是 $Q_A$-闭的。例如，2.2.2 节末的分析表明，自由手征多重态（相应地，自由扭手征多重态）的最低分量是手征（相应地，扭手征）局域算子。重要的是，由于 $[Q_A,Q_A]=[Q_B,Q_B]=0$，对任意局域算子 $O'$，$[Q_B,O']$ 平凡地是手征的，$[Q_A,O']$ 平凡地是扭手征的；我们把手征环定义为手征局域算子 $O$ 模去那些平凡地手征的局域算子，即手征环就是局域算子的 $Q_B$ 上同调。类似地，把扭手征环定义为局域算子的 $Q_A$ 上同调。

如下文将见，手征环与扭手征环自然具有分次交换代数的结构；事实上，在 3.1.3 节我们将看到它们还具有自然的（移位的）Poisson 括号。立刻可以看出：若 $\mathcal{T}$ 与 $\tilde{\mathcal{T}}$ 是互为镜像的理论，则在（例如）$\mathcal{T}$ 的手征环与 $\tilde{\mathcal{T}}$ 的扭手征环之间，必存在合适的代数同构。因此，若要检验 $\mathcal{T}$ 与 $\tilde{\mathcal{T}}$ 的镜像对称猜想，把 $\mathcal{T}$ 中的全部局域算子与 $\tilde{\mathcal{T}}$ 中的局域算子对应起来这一（可能极难的）任务，可以先通过（较易、必要但不充分的）匹配两者手征环与扭手征环的任务来检验。

下面说明手征环自然具有分次交换代数的结构；扭手征环的相应结果只需把 $Q_B$ 换成 $Q_A$ 即可得到。首先注意，局域算子的 $Q_B$ 上同调不依赖于局域算子的插入点。给定手征算子 $O$，一个简短的计算表明它在世界面上的平移 $\partial_z,\partial_{\bar z}$ 是 $Q_B$-恰当的（扭手征算子与 $Q_A$-恰当性同理），例如：
\[
\frac{i}{2}\partial_{\bar z}O=[P_{\bar z},O]=[[\bar Q_+,Q_+],O]=[Q_B,[Q_+,O]]\,.
\]
% [Tutor Note: 上式用到两步：其一，超对称代数中的反对易子 $\{Q_\alpha,\bar Q_\beta\}\propto P$（其中含 $P_{\bar z}$）把平移算子写成超荷的反对易子，故 $[P_{\bar z},O]$ 可写成超荷的双对易子；其二，用 Jacobi 恒等式把它整理为 $[Q_B,[Q_+,O]]$ 的形式。于是 $\partial_{\bar z}O$ 是 $Q_B$-恰当项，上同调类与插入点无关——这正是手征环能定义"乘积"的前提。注意 $[Q_+,O]$ 这种"超荷作用在算子上"的构造正是 §3 中拓扑下降（descent）与次级乘积的起点。]
% [Tutor Note: 物理图像：$Q_B$-闭算子给出 $B$ 模型的手征环，$Q_A$-闭算子给出 $A$ 模型的扭手征环；镜像对称把这两族环互换，因此比较环结构是检验镜像对偶的第一道（必要但不充分的）关口。]

% ==== 原文 PDF 第 19 页 ====
给定两个手征算子，它们的乘积仍是手征算子（扭手征情形同理）。乘积只需通过碰撞算子（即令它们取重合点）来定义。由拓扑不变性而来的位置无关性意味着：当两点重合时，乘积必须非奇异（至多差 $Q_B$-恰当项）。于是，取 $Q_B$ 上同调就把这些手征算子的完整算子乘积约化为一个普通乘积，把算子代数约化为通常意义上的代数！此外，手征算子平移的 $Q_B$-恰当性意味着，碰撞一组手征局域算子时顺序无关紧要：手征算子的代数乘积是结合的。由于总存在平凡局域算子 $1$（我们假定它不是 $Q_B$-恰当的），可见这是一个含幺的结合代数。最后，由于做乘积时没有优先方向，这个结合乘积还必须分次交换——交换手征算子时出现的符号只来自算子的费米子奇偶性。

\subsubsection{手征环与 Dolbeault 上同调}

先分析带 Calabi--Yau 靶空间 $\mathcal{X}$ 的二维 $\mathcal{N}=(2,2)$ $\sigma$ 模型在 $B$ 型扭变下的局域算子，即该理论的手征环。假定有 $N$ 个手征多重态 $\Phi^n$（及其共轭反手征多重态 $\bar\Phi^{\bar n}$），其轴 $R$ 荷与矢量 $R$ 荷均为零。利用手征场轴 $R$ 荷为零这一点，按 $B$ 型扭变自旋 $J_B=J+\frac12R_A$ 与上同调分次 $C_B=-R_V$ 重新组织这些场是直接了当的；我们把自旋 $J$、$R$ 荷、扭变自旋 $J_B$ 与上同调分次列于表 1。

\begin{table}[htbp]
\centering
\caption{自旋 $J$、轴与矢量 $R$ 荷 $(R_A,R_V)$、扭变自旋 $J_B=J+\frac12R_A$ 以及上同调分次 $C_B=-R_V$，对应 $R$ 荷为零的 $B$ 型扭变手征与反手征多重态。（原文第 19 页，Table 1）}
\begin{tabular}{lcccccc}
\toprule
 & $\phi$ & $\bar\phi$ & $\psi_\pm$ & $\bar\psi_\pm$ & $F$ & $\bar F$\\
\midrule
$J$ & 0 & 0 & $\pm\frac12$ & $\pm\frac12$ & 0 & 0\\
$(R_A,R_V)$ & $(0,0)$ & $(0,0)$ & $(\pm1,1)$ & $(\mp1,-1)$ & $(0,2)$ & $(0,-2)$\\
$J_B$ & 0 & 0 & $\pm1$ & 0 & 0 & 0\\
$C_B$ & 0 & 0 & $-1$ & 1 & $-2$ & 2\\
\bottomrule
\end{tabular}
\end{table}
% [Tutor Note: 表 1 的数值可直接由定义核对：$J_B=J+\frac12R_A$ 给出 $\psi$ 的 $J_B=\pm\frac12+\frac12(\pm1)=\pm1$、$\bar\psi$ 的 $J_B=\pm\frac12+\frac12(\mp1)=0$；$C_B=-R_V$ 给出 $\psi$ 的 $C_B=-1$、$\bar\psi$ 的 $C_B=+1$、$F$ 的 $C_B=-2$、$\bar F$ 的 $C_B=+2$。特别地 $\bar\psi$ 的轴荷写成 $(\mp1,-1)$ 正是为了让 $J_B=0$：镜像（手征↔反手征）把 $R_A$ 反号而保持 $R_V$。]

现在把这些 $B$ 型扭变的数据组织成一些自然的几何对象。首先，玻色子 $\phi,\bar\phi,\bar F$ 仍是世界面上的标量；我们继续把 $\phi,\bar\phi$ 看作映射 $\phi:\Sigma\to\mathcal{X}$（从世界面 $\Sigma$ 到 Calabi--Yau 靶空间 $\mathcal{X}$）的全纯/反全纯部分，并把辅助玻色子 $\bar F$ 对偶化为取值于全纯余切丛拉回的标量 $\bar F\in\phi^*T^*_{(0,1)}\mathcal{X}$。\footnote{严格来说，产生这个 $\phi^*T^*_{(1,0)}\mathcal{X}$ 取值的标量的是 $\bar F_n-\Gamma_{n\bar m\bar l}\bar\psi^{\bar m}_-\bar\psi^{\bar l}_+$。按记号滥用，我们把用凯勒度规对偶化后的这个移位场仍称为同一名字。对下文所述的世界面 2-形式 $F$ 也做同样的处理。}手征多重态中的费米子 $\psi^n_+,\psi^n_-$ 自然被重新组织为世界面上的 1-形式 $\rho^n=\frac{1}{2i}(-\psi^n_-dz+\psi^n_+d\bar z)$，\footnote{我们取世界面的微分形式 $dz,d\bar z$ 为费米子的。还发现给它们上同调次数 $1$ 比较方便：
\[
[C_B,dz]=dz\qquad [C_B,d\bar z]=d\bar z
\]
于是世界面外微分 $d_\Sigma=dz\partial_z+d\bar z\partial_{\bar z}$ 是上同调次数 $1$ 的（费米子）导子。由此，$\rho$ 自然是上同调次数 $0$ 的玻色型世界面 1-形式。}仍取值于全纯切丛的拉回，$\rho\in\Omega^1(\Sigma,\phi^*T^{(1,0)}\mathcal{X})$；而反手征
% [存疑: 本页（及下一页）对 $\phi^*T^*$ 的 $(1,0)$/$(0,1)$ 指标前后不一致：正文此处写 $\bar F\in\phi^*T^*_{(0,1)}\mathcal{X}$，而脚注 19 写 $\phi^*T^*_{(1,0)}\mathcal{X}$；第 20 页又把 $\bar\zeta$ 先写成 $\phi^*T^*_{(1,0)}\mathcal{X}$ 取值、后写成 $\phi^*T^*_{(0,1)}\mathcal{X}$ 取值。由指标位置判断（$\bar F_n$、$\bar\zeta_n$ 都带全纯下标 $n$，故应为全纯余切丛），似以 $(1,0)$ 为是，$(0,1)$ 疑为笔误。此处一律照抄原样，建议通稿时统一为 $\phi^*T^*_{(1,0)}\mathcal{X}$。]
% [原文疑误: 脚注 19 中 "The same holds to for the worldsheet 2-form F described below" 多出一个 "to"，应为 "holds for"。]
% [跨页接续: 原文第19页/第20页]
% ==== 原文 PDF 第 20 页 ====
多重态中的费米子 $\bar\psi^{\bar n}_\pm$ 成为世界面标量。我们发现，把这些费米子组织为（费米子、上同调次数 $1$）的 $\phi^*T^{(0,1)}\mathcal{X}$ 取值的标量 $\bar\psi^{\bar n}=\bar\psi^{\bar n}_++\bar\psi^{\bar n}_-\in\phi^*T^{(0,1)}\mathcal{X}$，以及（费米子、上同调次数 $1$）的 $\phi^*T^*_{(1,0)}\mathcal{X}$ 取值的标量 $\bar\zeta_n=g_{n\bar n}(\bar\psi^{\bar n}_+-\bar\psi^{\bar n}_-)$，比较方便。最后，玻色子 $F$ 自然成为（玻色子、上同调次数 $0$）取值于全纯切丛拉回的世界面 2-形式 $F=-\frac14F\,dz\,d\bar z\in\Omega^2(\Sigma,T^{(1,0)}\mathcal{X})$。有了这些场重定义，可以发现 $Q_B$ 的作用为
\begin{equation}\tag{2.30}
\begin{array}{l@{\quad}l@{\quad}l@{\quad}l}
Q_B\phi=0 & Q_B\rho=d_\Sigma\phi & Q_B F=d_\Sigma\rho & \\[4pt]
Q_B\bar\phi=\bar\psi & Q_B\bar\psi=0 & Q_B\bar\zeta=\bar F & Q_B\bar F=0
\end{array}
\end{equation}
其中 $d_\Sigma\phi\in\Omega^1(\Sigma,\phi^*T^{(1,0)}\mathcal{X})$ 是映射 $\phi:\Sigma\to\mathcal{X}$（的全纯部分）的（世界面）微分，而 $d_\Sigma\rho\in\Omega^2(\Sigma,\phi^*T^{(1,0)}\mathcal{X})$ 是 1-形式 $\rho$ 的（世界面）外微分 $d_\Sigma=dz\partial_z+d\bar z\partial_{\bar z}$。
% [Tutor Note: (2.30) 是 $B$ 型扭变的核心：$Q_B$ 的作用恰好是"把形式次数加一"的外微分——$\phi$ 是 0-形式而 $Q_B\phi=0$，$\rho$ 是 1-形式而 $Q_B\rho=d\phi$，$F$ 是 2-形式而 $Q_B F=d\rho$。核对时要同时检查形式次数与 $C_B$ 分次：$\rho$ 的 $C_B=0$、$F$ 的 $C_B=0$，而 $Q_B$ 不改变 $C_B$，只改变世界面形式次数。这一"超荷 = 外微分"的对应正是把上同调算子代数变成分次交换代数的机制。]

可以把 $B$ 型扭变中的局域算子构造为世界面 0-形式场的 $Q_B$-闭函数。具体说，就是玻色场 $\phi,\bar\phi$（它们参数化映射 $\phi:\Sigma\to\mathcal{X}$）、拉回全纯余切丛 $\phi^*T^*_{(1,0)}\mathcal{X}$ 的截面 $\bar F$，以及费米场 $\bar\psi$（取值于 $\phi^*T^{(0,1)}\mathcal{X}$）与 $\zeta$（取值于 $\phi^*T^*_{(0,1)}\mathcal{X}$）的函数。我们可以把这些玻色场与费米场的函数，等同于取值于全纯切丛外积的 $(0,\bullet)$-形式复形的截面 $\Omega^{(0,\bullet)}(\mathcal{X},\wedge^\bullet T^{(0,1)}\mathcal{X})$（沿 $\phi:\Sigma\to\mathcal{X}$ 拉回到世界面）。特别地，有如下对应：
% [原文疑误: 此处外积丛写作 $\wedge^\bullet T^{(0,1)}\mathcal{X}$，但同一句的正文说 "外积" 取自 "全纯切丛"（即 $T^{(1,0)}$），且 (2.32) 与紧随其后的讨论都用 $\wedge^\bullet T^{(1,0)}\mathcal{X}$，故 $T^{(0,1)}$ 应为 $T^{(1,0)}$ 之误。照抄原样，建议通稿时改为 $\wedge^\bullet T^{(1,0)}\mathcal{X}$。]
\begin{equation}\tag{2.31}
\bar\psi^{\bar n}\leftrightarrow d\bar\phi^{\bar n}\qquad\qquad \bar\zeta_n\leftrightarrow\partial_{\phi^n}
\end{equation}
% [Tutor Note: (2.31) 是 $B$ 型的"费米子—形式"字典：$\bar\psi^{\bar n}$ 对应 $(0,1)$-形式 $d\bar\phi^{\bar n}$，$\bar\zeta_n$ 对应全纯切向量 $\partial_{\phi^n}$。因此这些场的函数自然落在 $\Omega^{(0,\bullet)}(\mathcal{X},\wedge^\bullet T^{(1,0)}\mathcal{X})$ 中；后者由 Hochschild--Kostant--Rosenberg 定理同构于凝聚层导出范畴的 Hochschild 上同调，这就是 $B$ 模型的闭弦算子代数。]
% [背景补充: "多向量场"（polyvector field）指 $\wedge^\bullet T^{(1,0)}\mathcal{X}$ 的截面，即形如 $f(\phi)\,\partial_{\phi^{n_1}}\wedge\cdots\wedge\partial_{\phi^{n_k}}$ 的对象；它与微分形式通过全纯体积形式 $\Omega$ 对偶，是 $B$ 模型手征环的标准几何载体。]

由于运动方程 $\bar F=0$，我们无需考虑依赖 $\bar F$ 的函数：这意味着 Ward 恒等式 $\langle\bar F_n\dots\rangle=0$。（若有超势，则同样把 $\bar F_n$ 换成 $W_n$。）利用这一对应，超荷 $Q_B$ 被重新解释为上述复形上的 Dolbeault 微分：$Q_B\leftrightarrow\bar\partial_{\mathcal{X}}=d\bar\phi^{\bar n}\partial_{\bar\phi^{\bar n}}$。于是我们得出结论：$B$ 型扭变中的局域算子由向量丛 $\wedge^\bullet T^{(1,0)}\mathcal{X}$ 的 Dolbeault 上同调的元素标记：
\begin{equation}\tag{2.32}
\text{vector space of local operators:}\quad H^{(0,\bullet)}(\mathcal{X},\wedge^\bullet T^{(1,0)}\mathcal{X}):=H\big(\Omega^{(0,\bullet)}(\mathcal{X},\wedge^\bullet T^{(1,0)}\mathcal{X}),\bar\partial_{\mathcal{X}}\big).
\end{equation}
% [Tutor Note: 把辅助场 $\bar F$ 排除在外，是因为它的运动方程使其成为 $Q_B$-恰当的（等价地 $\langle\bar F_n\cdots\rangle$ 只给出接触项）；若有超势，$W_n$ 替代 $\bar F_n$ 成为同一复形里的元素，这正是 §3 中把 Landau--Ginzburg $B$ 模型算子代数识别为 $\bar\partial$-上同调（矩阵分解侧）的入口。]

正如预料，这些局域算子是 $\mathbb{Z}^2$-分次的：取值于 $\wedge^qT^{(1,0)}\mathcal{X}$ 的 $(0,p)$-形式，其上同调次数为（其矢量 $R$ 荷的负值）$q+p$，内部次数为（本质上\footnote{具有确定 $R_A$ 荷的算子（如上所定义）由原始的费米场 $\bar\psi^{\bar n}_\pm$ 构成。此处给出的内部次数仍是 $B$ 型扭变理论的对称性；它把 $-R_A$ 与旋转 $\bar\psi^{\bar m}_\pm\to\frac{1}{\sqrt2}(\bar\psi^{\bar m}_+\pm\bar\psi^{\bar m}_-)$ 结合起来。这个费米场空间的旋转不与通常的旋转对易（因为它混合不同自旋的场），但它确实与扭变后的旋转对易，从而是扭变理论的合法内部对称性。}是其轴 $R$ 荷的负值）$p-q$。
% [Tutor Note: 易错预警：内部次数 $p-q$ 只是"本质上"等于 $-R_A$——按脚注 21，具有确定 $R_A$ 荷的算子由原始的 $\bar\psi^{\bar n}_\pm$ 构成，而 $B$ 型扭变的内部对称性要把 $\bar\psi^{\bar n}_\pm$ 旋转为 $\frac{1}{\sqrt2}(\bar\psi^{\bar n}_+\pm\bar\psi^{\bar n}_-)$；该旋转混合不同自旋的场，故不与通常的旋转对易，却与扭变后的旋转对易。因此 $p-q$ 是扭变理论的内部次数，而不要与未扭变理论的 $R_A$ 荷直接等同。]

考虑仿射空间的情形：$\mathcal{X}=\mathbb{C}^N$。由 Poincaré 引理的全纯版本，每个 $\bar\partial_{\mathbb{C}^N}$-闭的 $(0,p)$-形式都是 $\bar\partial_{\mathbb{C}^N}$-恰当的。特别地，上同调群 $H^{(0,p)}(\mathbb{C}^N,\wedge^qT^{(1,0)}\mathbb{C}^N)$ 当 $p=0$ 时才不消失。$\bar\partial_{\mathbb{C}^N}$-闭的 $(0,0)$-形式就是 $\phi$ 的全纯函数，因此我们断言 $H^{(0,0)}(\mathbb{C}^N,\wedge^qT^{(1,0)}\mathbb{C}^N)$ 由全纯多向量场组成：
\begin{equation}\tag{2.33}
H^{(0,\bullet)}(\mathbb{C}^N,\wedge^\bullet T^{(1,0)}\mathbb{C}^N)=H^{(0,0)}(\mathbb{C}^N,\wedge^\bullet T^{(1,0)}\mathbb{C}^N)\cong\big\{f(\phi)\partial_{\phi^{n_1}}\wedge\dots\wedge\partial_{\phi^{n_k}}\big\}\cong\mathbb{C}[\phi^n,\bar\zeta_n].
\end{equation}
% [Tutor Note: (2.33) 是 $B$ 模型手征环的最简单例子：仿射空间上只剩全纯多向量场，且 $\bar\zeta_n$ 与 $\partial_{\phi^n}$ 的对应把它写成多项式环 $\mathbb{C}[\phi^n,\bar\zeta_n]$。若靶空间紧致（如 Calabi--Yau 流形），$H^{(0,p)}$ 一般不为零，其上同调类就是 $B$ 模型的关联函数所要积分的对象（见第 21 页的讨论）。]

% ==== 原文 PDF 第 21 页 ====
$Q_B$ 的不动点轨迹由世界面到靶空间的局部常值映射组成，即 $Q_B$ 不变的场组态满足 $d_\Sigma\phi=0$。这意味着上述局域算子的插入与其在世界面上的位置无关。此外，常值映射的空间就是 $\mathcal{X}$ 本身，且这些点状映射不接收任何非微扰的弦论修正。关于局域化及其在计算超对称可观测量（例如手征环元素的关联函数）中的作用，可参看附录 A 的简要回顾。粗略地说，$B$ 模型中关联函数的计算归结为把 Dolbeault 上同调类 $[\omega]\in H^{(0,\bullet)}(\mathcal{X},\wedge^\bullet T^{(1,0)}\mathcal{X})$ 在 $\mathcal{X}$ 上积分，这是一个经典的几何问题，不同于我们即将讨论的远为微妙的 $A$ 模型一侧。因此，$B$ 型扭变中局域算子的向量空间与微分形式的（经典）楔积一致——它不接收任何量子修正！

事实上，更细致的分析表明整个 $B$ 模型作用是 $Q_B$-恰当的。要看出 $B$ 模型只依赖靶流形的复结构则稍微困难一些，但这是正确的：$B$ 模型同时是世界面上的拓扑理论与靶空间中的“半拓扑”理论，依赖于与 $A$ 模型互补的那些几何形变。看清这一点的一种办法是：把上同调类 $\omega\in\Omega^{(0,q)}(\mathcal{X},\wedge^pT^{(1,0)}\mathcal{X})$ 的（某个代表元）在 $\mathcal{X}$ 上与 Calabi--Yau 的全纯体积形式 $\Omega$ 缩并，得到 Dolbeault 次数为 $(n-p,q)$ 的普通微分形式；只有当 $p=q=n$ 时，才能把这样的微分形式在 $\mathcal{X}$ 上积出消失的结果，而全纯体积形式依赖于复结构的选择。
% [存疑: 这里 "$(n-p,q)$-形式在实维数 $2n$ 的 $\mathcal{X}$ 上积分" 要求 $n-p+q=2n$，而原文条件 "$p=q=n$" 只给出 $(0,n)$-形式（只能在 $n$ 维闭链上积分），似不自洽；若原意是 $(2n-p,q)$，则 $p=q=n$ 恰给出顶形式 $(n,n)$，可以积分。判断为原文此处表述或指标有误，但无法确证，照抄原样并建议按 $(2n-p,q)$ 核对。]

\subsubsection{扭手征环与量子上同调}

下面转向二维 $\mathcal{N}=(2,2)$ Calabi--Yau $\sigma$ 模型的 $A$ 型扭变。我们仍取超势为零，并把它看作具有合适凯勒势的多个手征多重态的理论；同时取手征多重态的矢量 $R$ 荷与轴 $R$ 荷为零。相应的扭变自旋与上同调分次列于表 2。

\begin{table}[htbp]
\centering
\caption{扭变自旋 $J_A=J+\frac12R_V$ 与上同调分次 $C_A=-R_A$，对应 $R$ 荷为零的 $A$ 型扭变手征与反手征多重态。（原文第 21 页，Table 2）}
\begin{tabular}{lcccccc}
\toprule
 & $\phi$ & $\bar\phi$ & $\psi_-,\bar\psi_+$ & $\psi_+,\bar\psi_-$ & $F$ & $\bar F$\\
\midrule
$J_A$ & 0 & 0 & 0 & $1,-1$ & 1 & $-1$\\
$C_A$ & 0 & 0 & 1 & $-1$ & 0 & 0\\
\bottomrule
\end{tabular}
\end{table}
% [Tutor Note: 表 2 的数值可由定义核对：$\psi_-$ 有 $J=-\frac12$、$R_V=-1$，故 $J_A=-\frac12+\frac12(-1)=-1$，而表中把 $(\psi_-,\bar\psi_+)$ 合列并给出 $J_A=0$——因为 $\bar\psi_+$ 的 $J=+\frac12$、$R_V=-1$，给出 $J_A=+\frac12-\frac12=0$；同理 $(\psi_+,\bar\psi_-)$ 一列的两个值是 $\psi_+$ 的 $J_A=\frac12+\frac12=1$ 与 $\bar\psi_-$ 的 $J_A=-\frac12-\frac12=-1$。$C_A=-R_A$ 同样逐项吻合：$\psi_-$ 的 $R_A=-1$、$\bar\psi_+$ 的 $R_A=-1$（故合并列 $C_A=1$），$\psi_+$ 的 $R_A=1$、$\bar\psi_-$ 的 $R_A=1$（故 $C_A=-1$）。合并列的两个值分别对应未加横线与加横线的场，排序与原文一致。]

就分量场而言，我们有玻色子 $\phi,\bar\phi$，它们继续给出映射 $\phi:\Sigma\to\mathcal{X}$。与 $B$ 型扭变一样，两个费米子世界面旋量——这次是 $\psi_-$ 与 $\bar\psi_+$——成为（费米子、次数 $1$）的世界面标量，从几何上实现（拉回的）全纯切丛的截面 $\chi:=\psi_-\in\phi^*T^{(1,0)}\mathcal{X}$ 与反全纯切丛的截面 $\bar\chi:=\bar\psi_+\in\phi^*T^{(0,1)}\mathcal{X}$。另外两个费米子 $\psi_+$ 与 $\bar\psi_-$，以及辅助玻色子 $F,\bar F$，自然被等同于全纯/反全纯世界面 1-形式的分量：费米子成为 1-形式 $\bar\xi:=-\frac{1}{2i}\bar\psi_-dz$ 与 $\xi:=\frac{1}{2i}\psi_+d\bar z$（都是玻色子、次数 $0$），而玻色子成为 $\bar F=-\frac{1}{2i}\bar Fdz$ 与 $F=\frac{1}{2i}Fd\bar z$（费米子、次数 $1$）。由此，$Q_A$ 对这些场的作用为
\begin{equation}\tag{2.34}
\begin{array}{l@{\quad}l@{\quad}l@{\quad}l}
Q_A\phi=\chi & Q_A\chi=0 & Q_A\xi=\bar\partial_\Sigma\phi+F & Q_A F=\bar\partial_\Sigma\chi\\[4pt]
Q_A\bar\phi=\bar\chi & Q_A\bar\chi=0 & Q_A\bar\xi=\partial_\Sigma\bar\phi-\bar F & Q_A\bar F=-\partial_\Sigma\bar\chi
\end{array}
\end{equation}
% [Tutor Note: (2.34) 与 (2.30) 的对偶是本节的关键：$B$ 型扭变里 $\phi$ 是 0-形式、$Q_B\phi=0$；$A$ 型扭变里 $Q_A\phi=\chi\neq0$，超荷把标量送到切丛取值的费米子 $\chi$。因此在 $A$ 型扭变中单个手征多重态的最低分量不是 $Q_A$-闭的，$Q_A$-上同调须由 $\chi$、$\xi$ 等分量组合而成。物理上，$A$ 模型的关联函数由世界面瞬子（伪全纯曲线）贡献，因而依赖凯勒结构，与 $B$ 模型依赖复结构互补。]
% [Tutor Note: 易错预警：这里 $\bar\xi$、$\xi$ 与 $\bar F$、$F$ 各出现两次——一次指原本的超场分量，一次指其 1-形式化身，原文按脚注 19 的惯例以同一字母表示（"用凯勒度规对偶化后的移位场仍称为同一名字"）。抄录与核对时不要误以为 (2.34) 中有两套不同的场。]
% ==== 原文 PDF 第 22 页 ====
% [跨页接续: 原文第21页/22页，本句为 (2.34) 的说明，位于第 22 页首行]
% [边界说明: 本句归 s02e；s02f 自第 22 页第二段起]
其中 $\partial_\Sigma$、$\bar\partial_\Sigma$ 是世界面 Dolbeault 算子。
